\documentclass[11pt]{article}
\usepackage{ochamath}
\subjectcolor{2F5DA8}
\setsheet{線形代数・電気系院試補充}{最小多項式　演習}{https://university-math-crj.pages.dev/linear-algebra/minimal-polynomial/}
\namelinetrue
\begin{document}
\makesheethead{問題}
自作問題。本文の例題とは数値や設定が異なります。条件・途中式・検算を示してください。
\problem[最小と特性の違い]
$A=\operatorname{diag}(4,4,4,-2)$ の特性多項式と最小多項式を求めよ。
\workspace{45mm}
\problem[非自明なブロック]
$A=\operatorname{diag}(J_2(4),-2)$ の最小多項式を求め、対角化を判定せよ。
\workspace{45mm}
\newpage
\problem[整除性の証明]
$q(A)=0$ なら $m_A\mid q$ を証明せよ。
\workspace{45mm}
\problem[体の条件]
$A=\begin{pmatrix}0&-3\\3&0\end{pmatrix}$ の最小多項式を求め、実数上と複素数上で対角化を判定せよ。
\workspace{45mm}
\newpage
\problem[スペクトル射影]
$A^2=5A-6I$。固有値2,3への射影を $A$ で表し、$A^k$ を導け。
\workspace{45mm}
\problem[十分性の証明]
最小多項式が相異なる一次因子の積なら対角化可能であることを補間多項式で証明せよ。
\workspace{45mm}
\end{document}
